\section{Exact detecting constants for every endpoint multiplicity pattern}
\label{sec:analytic-orbital-detection-constants}

\begin{definition}[The Hermite orbital positivity radius]
For $m\ge2$ let
\[
 D_m=\operatorname{diag}(h_{m,1},\ldots,h_{m,m}),
\]
where the entries are all roots of the probabilists' Hermite polynomial
$\operatorname{He}_m$. Define
\[
 \phi_m(X)=\int_{U(m)}e^{i\operatorname{Tr}(D_mUXU^*)}\,dU,
                 \qquad X=X^*,
\]
and
\begin{equation}\label{eq:analytic-orbital-radius}
 \mathfrak r_m=
   \inf\{\lambda_{\max}(X)-\lambda_{\min}(X):\phi_m(X)<0\}.
\end{equation}
Put $\mathfrak r_1=+\infty$ and $h_m=\max_j|h_{m,j}|$.
The orbital measure is the concrete Haar measure on the compact
Hermite orbit, not a hypothesized arithmetic measure.
\end{definition}

\begin{lemma}[Strictness, bounds and the first two radii]
\label{lem:analytic-orbital-radius-properties}
For $m\ge2$,
\[
                  0<\frac\pi{h_m}\le\mathfrak r_m<\infty.
\]
For each $c<\mathfrak r_m$, $\phi_m$ has a strictly positive minimum
on the compact set of diagonal matrices with entries in $[0,c]$.
Moreover
\begin{equation}\label{eq:analytic-first-orbital-radii}
                   \mathfrak r_2=\pi,\qquad
                   \mathfrak r_3=\frac{2\pi}{\sqrt3}.
\end{equation}
\end{lemma}

\begin{proof}
Finiteness follows from Lemma~\ref{lem:analytic-negative-orbital}.
The trace of $D_m$ is zero, so scalar translation of $X$ does not
change $\phi_m$. For small spectral diameter, subtracting a scalar
makes $\|X\|$ uniformly small. Since $\phi_m(0)=1$, this also proves
that the radius is positive.

We first rule out interior zero values in a nonnegative region.
On the Euclidean space of Hermitian matrices, differentiating under
the compact integral gives the exact elliptic equation
\begin{equation}\label{eq:analytic-orbital-Helmholtz}
                         \Delta\phi_m=-\|D_m\|^2\phi_m.
\end{equation}
The function is real analytic and real valued. If it were nonnegative
near a zero $X_0$, let $P$ be the first nonzero homogeneous term in
its Taylor series there. Such a term exists since $\phi_m$ is not
identically zero. Nonnegativity makes $P$ nonnegative in every real
direction. Equation~\eqref{eq:analytic-orbital-Helmholtz} forces its
lowest-degree part to satisfy $\Delta P=0$. The mean-value property
of harmonic polynomials gives average $P= P(0)=0$ on a sphere, so
nonnegative $P$ must vanish identically, a contradiction. Thus every
zero is approached by negative values. For $c<\mathfrak r_m$, a
neighborhood of every matrix of spectral diameter at most $c$ lies
inside the region with no negative values. No zero can occur there,
and compactness gives the stated strictly positive minimum.

To obtain the explicit lower bound, fix a diagonal spectrum $L$ of
diameter less than $\pi/h_m$, and shift its minimum to zero.
For small $\delta>0$, the finite alphabet
$e^{i\delta h_{m,j}}$ has conjugate pairs of angles at most
$\delta h_m$. The quadratic sector proof in NN-A020 makes it
$PF_r$ whenever $(r+1)\delta h_m<\pi$. Choose the partition
$\mu_i=\lfloor L_i/\delta\rfloor$, in descending order.
Its width satisfies this inequality for sufficiently small $\delta$.
Its Schur value is nonnegative. The exact normalized Weyl--HCIZ
identity \eqref{eq:analytic-Schur-orbital-limit} then gives
$\phi_m(L)\ge0$. Hence no negative value has diameter less than
$\pi/h_m$, proving the bound.

For $m=2$, the HCIZ determinant with Hermite eigenvalues $-1,1$
gives
\[
                \phi_2(X)=\frac{\sin(\lambda_1-\lambda_2)}
                                      {\lambda_1-\lambda_2},
\]
with its continuous value at equality. Its first sign change gives
$\mathfrak r_2=\pi$. For $m=3$, the eigenvalues are
$-\sqrt3,0,\sqrt3$. Factoring the exponential Vandermonde in the
same exact formula gives
\[
 \phi_3(X)=\prod_{1\le i<j\le3}
            \operatorname{sinc}\left(
                          \frac{\sqrt3}{2}(\lambda_i-\lambda_j)\right),
                    \qquad\operatorname{sinc}(u)=\frac{\sin u}{u}.
\]
All factors are positive when the diameter is less than
$2\pi/\sqrt3$. For eigenvalues $0,L/2,L$ with
$2\pi/\sqrt3<L<4\pi/\sqrt3$, the largest-difference factor is
negative and the other two are positive. This proves the second
formula in \eqref{eq:analytic-first-orbital-radii}.
\end{proof}

\begin{lemma}[Uniform finite-cluster positivity]
\label{lem:analytic-uniform-cluster-PF}
Let $m$ be fixed and let a conjugation-stable finite alphabet have
entries $b_\varepsilon e^{\varepsilon a_j(\varepsilon)}$, where
$b_\varepsilon>0$ and
$a_j(\varepsilon)=i\kappa h_{m,j}+O(\varepsilon)$ with fixed
$\kappa>0$. For each $c<\mathfrak r_m/\kappa$, its coefficient
sequence is $PF_{r(\varepsilon)}$ for all sufficiently small
$\varepsilon$, uniformly over integers $r(\varepsilon)\varepsilon\le c$.
The assertion includes every partition width in that range, with no
bound on the degree of the Toeplitz minor.
\end{lemma}

\begin{proof}
The positive common scale does not affect signs. For every partition
$\mu$ of at most $m$ rows and width at most $r$, the exact formula is
\[
 \frac{s_\mu(e^{\varepsilon a})}{s_\mu(1^m)}
       =\frac{\mathcal I(A_\varepsilon,
                       \operatorname{diag}(\varepsilon(\mu+\rho)))}
              {\mathcal I(A_\varepsilon,
                       \operatorname{diag}(\varepsilon\rho))},
            \qquad A_\varepsilon=\operatorname{diag}(a_j),
            \quad \rho=(m-1,\ldots,0).
\]
The second argument in the numerator has entries in
$[0,c+(m-1)\varepsilon]$. Choose $c'$ strictly between $c$ and
$\mathfrak r_m/\kappa$. The preceding lemma gives a positive
minimum for $\phi_m(\kappa X)$ over entries in $[0,c']$.
The Haar integral is uniformly continuous on these compact matrix
sets, while $A_\varepsilon=i\kappa D_m+O(\varepsilon)$.
The numerator is therefore uniformly close to a positive real value,
and the denominator tends to one. The character ratio itself is real,
by conjugation stability. It follows that every such ordinary Schur
value is strictly positive, uniformly in this full partition range.
Partitions with more than $m$ rows have value zero.
Every skew Schur value of width at most $r$ is a nonnegative integer
combination of ordinary Schur values of that width, by the exact
restriction-multiplicity argument in NN-A021. Thus all Toeplitz minors
of order at most $r$ are nonnegative. This proof retains the full
$O(\varepsilon)$ perturbation and the possible coincident limiting
scaled row lengths in the compact set.
\end{proof}

\begin{lemma}[Cluster isolation with earlier multiple endpoints]
\label{lem:analytic-coalescent-prefix-isolation}
The negative-witness conclusion of
Lemma~\ref{lem:analytic-cluster-Schur-witness} remains valid when its
finite prefix is a conjugation-stable multiset of possibly complex
parameters $\zeta_1(\varepsilon),\ldots,\zeta_p(\varepsilon)$ with
positive real limits strictly larger than $b^0$, even when some limits
coincide, provided the parameters are distinct for $\varepsilon>0$ and
\[
                 |\zeta_i-\zeta_j|\ge c\varepsilon
                         \quad(i\ne j)
\]
for a fixed $c>0$. Real prefix entries are positive for small
$\varepsilon$, so $P=\prod_i\zeta_i$ is positive. For any strictly
negative orbital value at spectral diameter $d$, the witness width
can be chosen with
\[
                   \limsup_{\varepsilon\downarrow0}
                               \varepsilon\lambda_1\le d+\delta
\]
for any prescribed $\delta>0$.
\end{lemma}

\begin{proof}
By continuity, first perturb the negative orbital spectrum to have
distinct eigenvalues while preserving strict negativity and increasing
its diameter by less than $\delta/4$. Shift that spectrum to have minimum in $(0,\delta/4)$,
and choose the bounding first-row length constant $C$ less than
$d+\delta$ while keeping a fixed positive gap above its maximum.
Choose the cluster row lengths and the prefix rows of length
$n=\lceil C/\varepsilon\rceil$ exactly as in NN-A028. The gaps
$\mu_m\ge\eta/\varepsilon$ and
$n-\mu_1\ge\eta/\varepsilon$ still hold for some fixed $\eta>0$.
Strict negativity is preserved by the scalar shift because the
limiting Hermite diagonal is traceless.

Here are the changes to the complete determinant estimates.
All $q=p+m$ retained poles now have pairwise separation at least a
constant times $\varepsilon$; the prefix-to-cluster separation remains
bounded below by a positive constant. Thus every partial-fraction
coefficient and the fixed-circle analytic remainder are
$O(\varepsilon^{-(q-1)})$, rather than
$O(\varepsilon^{-(m-1)})$. In a term with remainder columns, the same
sorted-log-modulus argument applies. The entire prefix is assigned
exponent $n$ in the maximal full weight. Put
\[
                  W=P^n b^{|\mu|};
\]
the maximal full weight is $W e^{O(\varepsilon)}$.
The positivity of $P$ makes this the same positive reference scale
as before. Replacing one retained modulus by the smaller tail modulus
still costs a fixed factor to the power $\mu_m$. Consequently the
total tail error is
\[
                     O(W\varepsilon^{-K}q_1^{\mu_m})
\]
for fixed $K$ and $q_1<1$; only the harmless power $K$ has changed.

In the Weyl numerator the leading prefix block is still exactly
$P^{n+m}\Delta(\zeta_1,\ldots,\zeta_p)$. It cancels its entire
prefix Vandermonde, regardless of how small that Vandermonde is.
The leading cross factor is
$\prod_{i,j}(1-\xi_j/\zeta_i)^{-1}$; it is real and tends to a
positive number. Every wrong assignment still exchanges a larger
prefix modulus with a smaller cluster modulus between rows separated
by at least $n-\mu_1$. Its exponential loss is therefore
$q_2^{n-\mu_1}$ with $q_2<1$. Dividing by the full Vandermonde now
costs at most $O(\varepsilon^{-q(q-1)/2})$, by the stated separation.
The wrong-prefix error is again
\[
                      O(W\varepsilon^{-K'}q_2^{n-\mu_1}).
\]
Both errors are smaller than any power of $\varepsilon$ times $W$.
The leading cluster Schur value, divided by its positive dimension
and $b^{|\mu|}$, tends to the fixed strictly negative orbital value.
Its sign is therefore unchanged. The dual-Jacobi--Trudi minor has
order $\lambda_1\le n$, including the empty-prefix case, and fits
in window $0,\ldots,\lambda_1+q-1$. Taking $\varepsilon\to0$
gives the asserted width bound.
\end{proof}

\begin{theorem}[NN-A029: exact orbital formulas for detecting order, window and rank]
\label{thm:analytic-exact-orbital-detection}
Assume A and put $p_0=F_0'(0)>0$. For each positive endpoint zero
$x$ of multiplicity $m_x$, put $\beta_x=1/(p_0x^2)$. Define
\begin{align}
 \mathcal L_F&=\min_{x:m_x\ge2}
                       \frac{x\mathfrak r_{m_x}}{2\sqrt2},
                          \label{eq:analytic-exact-LF}\\
 \mathcal L_U&=\min_{x:m_x\ge2}
                       \frac{x(1-\beta_x)\mathfrak r_{m_x}}{2\sqrt2},
                          \label{eq:analytic-exact-LU}
\end{align}
where a minimum over the empty set means $+\infty$. If there is any
multiple zero, these positive finite minima are attained. Then, in
the extended positive real line,
\begin{equation}\label{eq:analytic-exact-orbital-limits}
 \lim_{a\downarrow0}\sqrt a\,m_*(-a)
 =\lim_{a\downarrow0}\sqrt a\,R_*(-a)=\mathcal L_F,
       \qquad
 \lim_{a\downarrow0}\sqrt a\,d_*(-a)=\mathcal L_U.
\end{equation}
These formulas contain no unstated endpoint multiplicity or simplicity
hypothesis. They express the constants in terms of the actual endpoint
multiplicities and explicit finite-dimensional orbital integrals;
they do not determine those multiplicities.

A double endpoint contributes $\pi x/(2\sqrt2)$ to
\eqref{eq:analytic-exact-LF}; a triple endpoint contributes
$\pi x/\sqrt6$. The corresponding unitary contributions are
multiplied by $1-\beta_x$. Higher-multiplicity contributions are
specified exactly by \eqref{eq:analytic-orbital-radius}, without
asserting a closed formula for every $\mathfrak r_m$.
\end{theorem}

\begin{proof}
The large-radius multiplicity estimate in NN-A022 gives
$h_{m_x}/x\to0$ along positive zeros tending to infinity.
Since $\mathfrak r_m\ge\pi/h_m$, the quantities
$x\mathfrak r_{m_x}$ tend to infinity along multiple zeros escaping
to infinity. Also $1-\beta_x\to1$. Thus each nonempty minimum is
positive, finite and attained among a finite set of endpoint zeros.

Write $\varepsilon=\sqrt a$. The local parameters near an endpoint
$x$ have the expansion, after removing their positive geometric mean,
\[
 z_j(-a)^{-2}=b_x(a)e^{\varepsilon a_{x,j}(\varepsilon)},\qquad
 a_{x,j}(\varepsilon)=-\frac{2i\sqrt2}{x}h_{m_x,j}
                                             +O(\varepsilon).
\]
Changing the sign of the Hermite diagonal only permutes it. At a
simple endpoint the parameter stays positive and real. Uniform
$\ell^1$ tails, compact zero convergence and positivity of the endpoint
products were proved in NN-A027--A028.

Fix any finite $c<\mathcal L_F$. Choose a large circle, avoiding
endpoint zeros, so that the uniform high-radius angular estimate in
NN-A022 gives, for all zeros outside it and all sufficiently small $a$,
\[
                      |\arg(z_j(-a)^{-2})|/\varepsilon
                                      <\frac{\pi}{2(c+1)}.
\]
Every high conjugate quadratic factor is consequently $PF_r$ when
$r=\lfloor c/\varepsilon\rfloor$, for small $\varepsilon$; positive
linear factors are $PF_\infty$. Inside the circle there are finitely
many endpoint clusters. For each multiple cluster,
$c<x\mathfrak r_{m_x}/(2\sqrt2)$, so
Lemma~\ref{lem:analytic-uniform-cluster-PF} gives $PF_r$ uniformly
for its actual perturbed factors. Simple clusters are positive linear
factors. Multiplication and locally uniform convergence now give
$PF_r$ for the full entire coefficient sequence.

In fact every allowed minor of order at most $r$ is strictly positive.
There are infinitely many high factors with constant coefficient one
and positive first coefficient, each $PF_r$. In a Cauchy--Binet path,
move one index at a time one step to its required target without
colliding with another index; each such elementary factor minor is
its positive first coefficient. A finite number of high factors
realizes this path, and all remaining factors may contribute their
positive diagonal minors. Every other term is nonnegative. This is
the explicit strict-path argument used in NN-A020, and it applies
equally to the conjugate quadratic factors here. Thus
$m_*(-a)>r$, and $R_*(-a)\ge m_*(-a)-1$. Since $c<\mathcal L_F$
was arbitrary, this proves the lower limits, also when
$\mathcal L_F=+\infty$.

For the upper limit when a multiple endpoint exists, choose any such
$x$ and a negative orbital value of the scaled Hermite diagonal
$D=(2\sqrt2/x)D_{m_x}$ with spectral diameter less than
$x\mathfrak r_{m_x}/(2\sqrt2)+\delta$.
Its existence follows directly from the infimum defining
$\mathfrak r_{m_x}$. All endpoint parameters below $x$ form a finite
prefix. They need not stay real: earlier multiple endpoints may split
into conjugate clusters. Their moduli, however, remain uniformly
larger than that of the chosen cluster; all limits are positive real.
Within each earlier cluster the Hermite roots give separation at least
a constant times $\varepsilon$, and between distinct endpoint clusters
the separation stays bounded below. Their product is positive.
Hence Lemma~\ref{lem:analytic-coalescent-prefix-isolation} applies.
It gives negative minors of order $n$ and window at most $n+O_x(1)$
with
\[
 \limsup_{a\downarrow0}\sqrt a\,n
                    \le\frac{x\mathfrak r_{m_x}}{2\sqrt2}+2\delta.
\]
Choose a minimizing endpoint and let $\delta\downarrow0$. Together
with the lower limits this proves both limits involving $\mathcal L_F$.

For the unitary statement use the actual translated alphabet
\[
                \widetilde\beta_j(a)=\frac1{p_a z_j(-a)^2-1},
                                \qquad p_a=F_{-a}'(0).
\]
At endpoint $x$ its scaled Hermite diagonal is
\[
                    \frac{2\sqrt2}{x(1-\beta_x)}D_{m_x}.
\]
This follows from the full local expansion in NN-A027; it retains the
$O(\varepsilon)$ perturbation after scaling. The positive endpoint
ordering, the finite coalescent-prefix separation and the uniform
$\ell^1$ tail remain valid. The high-radius angular tail remains
uniformly negligible after division by $\varepsilon$, since
$|\beta_j(a)|\le C/|z_j(-a)|^2$ and
$|\arg(\beta/(1-\beta))|\le|\arg\beta|/(1-|\beta|)$.
The same lower and upper arguments therefore give the constant
\eqref{eq:analytic-exact-LU} for the least negative-minor order of
the translated coefficients. NN-A021 identifies that order exactly
with the least failing unitary rank. Its positive constant
$E_{-a}(-1)$ does not change any minor sign. This proves the remaining
limit. The double and triple values follow from
\eqref{eq:analytic-first-orbital-radii}.
\end{proof}

\paragraph{Geometric and representation meaning.}
For each endpoint local Artin length $m$, the leading heat unfolding
specifies the actual full-flag coadjoint orbit of the Hermite diagonal.
Its Fourier transform has the exact finite-dimensional positivity
radius \eqref{eq:analytic-orbital-radius}. The normalized Weyl--HCIZ
identity controls every integral highest-weight partition in a whole
scaled-width compact set, while its first negative orbital value gives
a negative Schur coefficient after the proved coalescent-prefix and
infinite-tail isolation. These are actual exterior-power tests, with
the representation-ring formalism in SGA~6, Expos\'e~V, Example~2.2(b)
and Lemma~2.2.1. The minima in
\eqref{eq:analytic-exact-LF}--\eqref{eq:analytic-exact-LU} are the exact
asymptotic boundaries of the corresponding analytic positive cones.
The primary HCIZ normalization is Zuber2018 equations (6), (11)--(12),
pp.~311--313, cited in NN-A028. No arithmetic Frobenius object,
integrality of the Hermite symplectic form, or complexity-class
separation is inferred.
